
\noindent Para asegurar un avance controlado, auditable y evitar cuellos de botella entre los equipos, la documentación se ha desglosado en entregables intermedios y finales. Esto permite validar fases críticas antes de invertir horas de diseño en la siguiente etapa. Los entregables se dividen por categorías funcionales:

\section{Requisitos y Arquitectura de Sistema (Líderes)}
\noindent Documentos fundamentales que establecen las reglas del juego antes de comenzar cualquier diseño electrónico.

\begin{table}[htbp]
\centering
\renewcommand{\arraystretch}{1.5}
\begin{tabularx}{\textwidth}{>{\hsize=0.7\hsize}X >{\hsize=0.3\hsize}X >{\hsize=2.0\hsize}X}
\toprule
\textbf{Entregable} & \textbf{Formato} & \textbf{Descripción} \\ \midrule
\textbf{Auditoria proyecto anterior} & PDF & Revisión crítica del proyecto del equipo de años anteriores, valorar que parte del proyecto es de utilidad y que no es necesario. \\
\textbf{System Requirements Document (SRD)} & PDF & Requisitos de alto nivel: voltajes, corrientes límite, parámetros orbitales y modos de fallo. Base inamovible del diseño. \\ 
\textbf{Power Budget Master} & Excel & Hoja de cálculo con los consumos pico y medio de todos los subsistemas. \textbf{Requisito previo estricto} para el Trade-Off. \\ 
\textbf{Mechanical ICD (Restricciones Físicas)} & DXF / PDF & Definición del volumen de la PCB, posiciones de montaje y altura máxima (aportado por los agentes dobles de Estructuras). \\ \bottomrule
\end{tabularx}
\end{table}
\clearpage
\section{Diseño Electrónico Hardware (Equipo Asignatura)}
\noindent Documentación que justifica y refleja el diseño físico, desde la elección del primer chip hasta el plano eléctrico final.

\begin{table}[H]
\centering
\renewcommand{\arraystretch}{1.5}
\begin{tabularx}{\textwidth}{>{\hsize=0.7\hsize}X >{\hsize=0.3\hsize}X >{\hsize=2.0\hsize}X}
\toprule
\textbf{Entregable} & \textbf{Formato} & \textbf{Descripción} \\ \midrule
\textbf{Trade-Off Report} & Notion / PDF & Memoria justificativa de la selección de reguladores (ej. LTC3130), sensores y protecciones frente a alternativas comerciales. \\ 
\textbf{Bill of Materials (BOM) Preliminar} & Excel & Listado inicial de componentes con números de parte (MPN), encapsulados, disponibilidad en stock y precio para validación técnica. \\ 
\textbf{PSIM Simulation Report (v1.0 Ideal)} & PDF & Validación virtual rápida de la topología matemática de potencia usando componentes y medidores ideales. \\ 
\textbf{PSIM Simulation Report (v2.0 Real)} & PDF & Validación avanzada con modelos no ideales, integrando picos de arranque, actuación de Load Switches y transitorios térmicos. \\ 
\textbf{Schematics Package} & PDF / Netlist & Planos eléctricos finales en ECAD (KiCad/Altium) organizados jerárquicamente por bloques. Incluye la auditoría ERC (Electrical Rules Check). \\ \bottomrule
\end{tabularx}
\end{table}

\section{Diseño Lógico e Integración (Equipo SART)}
\noindent Documentos que definen el "cerebro" del sistema y cómo el EPS se comunica con el Ordenador de a Bordo (OBC).

\begin{table}[H]
\centering
\renewcommand{\arraystretch}{1.5}
\begin{tabularx}{\textwidth}{>{\hsize=0.7\hsize}X >{\hsize=0.3\hsize}X >{\hsize=2.0\hsize}X}
\toprule
\textbf{Entregable} & \textbf{Formato} & \textbf{Descripción} \\ \midrule
\textbf{State Machine Diagram} & Draw.io / PDF & Diagrama de flujo visual que mapea los estados del EPS (Nominal, Eclipse, Safe Mode, Critical Fault) y sus condiciones lógicas de transición. \\ 
\textbf{Interface Control Document (ICD)} & Excel & Documento de conexión absoluta. Contiene el mapeo exacto de pines (Pinout), líneas ADC/GPIO, frecuencias I2C/SPI y matriz de comandos OBC-EPS. \\ 
\textbf{CDR Acta (Critical Design Review)} & PDF / Notion & Acta formal de la reunión donde el equipo SART y los Líderes auditan los esquemas de la Asignatura antes de darlos por válidos, registrando fallos y correcciones. \\ \bottomrule
\end{tabularx}
\end{table}